% !TEX program = pdflatex
% Paper 2 (Molecular Informatics short / journal-length version)
% Condensed main text (~6 pp). Material removed here lives in the
% Supporting Information (paper2_SI.tex) and in the full-length draft
% (paper2_orthogonality_screening.tex). Compile twice for cross-refs.

\documentclass[11pt,a4paper]{article}

\usepackage[T1]{fontenc}
\usepackage[utf8]{inputenc}
\usepackage{lmodern}
\usepackage[margin=1in]{geometry}
\usepackage{amsmath,amssymb,amsthm}
\usepackage{mathtools}
\usepackage{booktabs}
\usepackage{array}
\usepackage{graphicx}
\graphicspath{{../figures/}}
\usepackage{hyperref}
\hypersetup{colorlinks=true,linkcolor=blue!50!black,citecolor=blue!50!black,urlcolor=blue!50!black}
\usepackage{xcolor}
\usepackage{enumitem}
\setlist[itemize]{topsep=2pt,itemsep=2pt}

\newtheorem{theorem}{Theorem}[section]
\theoremstyle{definition}
\newtheorem{remark}[theorem]{Remark}

\newcommand{\bid}{\mathcal{F}_{\mathrm{BID}}}

\title{\bfseries Orthogonality Screening of Topological Indices for QSPR Modeling:\\ A Structural and Empirical Redundancy Analysis}

\author{
  Ganesh Shiwakoti\thanks{Department of Mathematics and Statistics, Florida Atlantic University, USA. \texttt{ganshiwakoti@gmail.com}} \and
  Advik Natarajan\thanks{Department of Mathematics, Loyola College, Chennai, India.} \and
  Parik Chalise\thanks{Department of Applied Mathematics and Statistics, Johns Hopkins University, Baltimore, USA.} \and
  Michael Arockiaraj\footnotemark[2]
}

\date{Draft of \today}

\begin{document}
\maketitle

% ----- Graphical abstract / TOC entry (Molecular Informatics) -----
\begin{center}
\includegraphics[width=0.82\linewidth]{fig_graphical_abstract.pdf}
\end{center}
\noindent\textbf{TOC entry.}\enspace Hundreds of topological indices
exist, and the standard pairwise-correlation test for whether a new
one carries ``independent'' information is too permissive. Combining
pairwise correlation with partial correlation against the QSPR
target, only $0$--$4$ of $20$ candidate indices carry independent
signal across four MoleculeNet datasets even though $15$--$17$ clear
the pairwise screen; structurally, the entire bond-incident-degree
family is confined to a ten-dimensional space. Orthogonality is
necessary but not sufficient.

\begin{abstract}
The chemical-graph-theory literature contains hundreds of scalar
topological indices, and the usual test for whether a newly
proposed index carries independent information --- a pairwise
correlation below $0.95$ against a classical baseline --- is
empirically too permissive. We present an open-source screening
pipeline that combines pairwise correlation, principal-component
effective rank, variance inflation, and partial correlation with
the modelling target into a single pre-flight check, complemented
by a structural result (the \emph{Edge-Degree-Pair Basis}) showing
that every bond-incident-degree index on hydrogen-suppressed
molecular graphs of maximum degree four lies in a vector space of
dimension at most ten. Applied to a three-parameter parametric
index family and to twenty alternative candidate indices across
four MoleculeNet benchmarks, the screen shows that most candidates
clear the pairwise test yet add little independent signal once
partial correlation is also required; a cross-validated
machine-learning benchmark confirms that the pruned index set
preserves predictive performance at several-fold fewer features.
We conclude that orthogonality is necessary but not sufficient for
a new index to be useful, and provide the pipeline as a reusable
standard --- noting that the diagnostics are linear and that the
screen is a pre-flight check, not a hard feature-selection filter.

\medskip
\noindent\textbf{Keywords.} Topological index, orthogonality
screening, redundancy analysis, QSPR, principal component analysis,
partial correlation, bond-incident-degree (BID) family.

\medskip
\noindent\textbf{Code and data.}
\url{https://github.com/Singati2/topological-index-orthogonality}
\end{abstract}

\section{Introduction}

The vocabulary of scalar topological indices used in QSPR/QSAR
modelling has grown over six decades from a handful of
proposals~\cite{Wiener1947, Randic1975, Gutman1972} to several
hundred, and the resulting proliferation has been criticised within
mathematical chemistry itself~\cite{Gutman2013Degree}. Many indices
are interrelated by exact identity, and on a fixed dataset the
empirical Pearson correlations between proposed indices routinely
exceed $|r| = 0.95$ for a large fraction of pairs.

This raises a recurring methodological question: \emph{what must a
newly proposed scalar topological index demonstrate to count as a
meaningful contribution to QSPR modelling?} Three weak
justifications recur --- a strong correlation with a property on a
small dataset, a closed form on standard graph classes, and
parametric generality that subsumes a classical index at one
parameter value --- yet none establishes \emph{independent}
predictive value beyond an existing baseline. We propose an
explicit, reusable pre-flight check for the ``does it carry
independent information'' question.

Our contributions are: (1)~an open-source orthogonality-screening
pipeline combining a pairwise-correlation redundancy threshold,
principal-component effective rank, variance inflation, and partial
correlation with the modelling target; (2)~replication across four
MoleculeNet benchmarks --- three regression tasks (ESOL, FreeSolv,
Lipophilicity) and one classification task (BBBP); (3)~a structural
result, the \emph{Edge-Degree-Pair Basis}, bounding the
bond-incident-degree (BID) family within a ten-dimensional space on
molecular graphs of maximum degree four; (4)~two case studies --- a
three-parameter parametric index family~\cite{Paper1_Wuzi} screened
over $300$ grid points and twenty alternative-family candidate
indices; and (5)~the explicit positioning of orthogonality as
\emph{necessary but not sufficient} for QSPR usefulness.

\section{Methods}

\subsection{Datasets and baseline}

We screen against a $30$-index classical baseline following
practitioner-standard cheminformatics
practice~\cite{TodeschiniConsonniHandbook}: degree-based,
distance-based, spectral, and other-family indices in the
proportions in which they are used in published QSPR work
(eighteen of the thirty are degree-based BID indices). Indices are
computed on hydrogen-suppressed molecular graphs (largest connected
fragment) with NetworkX, RDKit~\cite{RDKit}, and
scikit-learn~\cite{Pedregosa2011sklearn}; each admits a
$\le 20$-line implementation. The four benchmarks are summarised in
Table~\ref{tab:datasets}; the full $30$-index list is given in the
Supporting Information.

\begin{table}[h]
\centering
\caption{QSPR/QSAR benchmark datasets used in this study; $n$ counts molecules surviving RDKit parsing (largest connected fragment, $\ge 2$ heavy atoms).}
\label{tab:datasets}
\begin{tabular}{lllrl}
\toprule
Name & Target property & Task & $n$ & Source\\
\midrule
ESOL          & $\log S$ (aqueous solubility, mol/L)         & regression     & $1\,127$ & \cite{Delaney2004ESOL} via \cite{MoleculeNet}\\
FreeSolv      & $\Delta G_{\mathrm{hyd}}$ (hydration free energy, kcal/mol) & regression & $639$    & \cite{Mobley2014FreeSolv} via \cite{MoleculeNet}\\
Lipophilicity & $\log D$ at pH 7.4                            & regression     & $4\,200$ & \cite{LipophilicityChEMBL} via \cite{MoleculeNet}\\
BBBP          & blood-brain barrier penetration               & classification & $2\,039$ & \cite{Martins2012BBBP} via \cite{MoleculeNet}\\
\bottomrule
\end{tabular}
\end{table}

\subsection{The Edge-Degree-Pair Basis}

Every BID index is a linear functional of edge-degree-pair counts,
which confines the entire family to a fixed finite-dimensional
space.

\begin{theorem}[Edge-Degree-Pair Basis]\label{thm:bid}
Let $\mathcal{G}_\Delta$ denote the class of finite simple graphs
with maximum degree at most $\Delta$. For any BID index
$I_f\in\bid$ and any $G\in\mathcal{G}_\Delta$,
\begin{equation}\label{eq:bid_linear}
I_f(G)=\sum_{1\le i\le j\le\Delta} f(i,j)\cdot m_{ij}(G),
\end{equation}
where $m_{ij}(G)=\bigl|\{uv\in E(G):\{d_u,d_v\}=\{i,j\}\}\bigr|$.
Restricted to $\mathcal{G}_4$, the vector space spanned by all BID
indices has dimension at most $10$.
\end{theorem}

The proof --- a partition of the edge set by endpoint-degree pair
followed by the dimension count $N_\Delta = \Delta(\Delta+1)/2$
(so $N_4 = 10$) --- is given in the Supporting Information. The
observation that BID indices are linear in the $m_{ij}$ counts is
standard~\cite{BorovicaninDasFurtulaGutman2017Zagreb}, and the
finite-dimensional linear-algebraic view of vertex-degree-based
indices has recently been developed in general by
Rada~\cite{Rada2026LinearVDB}. The consequence is operational: any
new BID index proposed for molecular chemistry ($\Delta\le 4$) must
justify its independence from the existing baseline within a
ten-dimensional subspace.

\subsection{Screening protocol}

For a candidate index $z$ and baseline $X=\{X_b\}$ on a dataset with
target $y$, we compute four diagnostics. (i)~\emph{Pairwise
redundancy}: the maximum absolute Pearson correlation
$\max_b |r(z, X_b)|$; a value below the conventional threshold
$0.95$ is the usual ``independence'' test. (ii)~\emph{Effective
rank}: the number of principal components of $X$ explaining $95\%$
and $99\%$ of variance~\cite{WoldEsbensenGeladi1987PCA}, measuring
baseline collinearity. (iii)~\emph{Variance inflation}: the VIF of
$z$ regressed on $X$~\cite{OBrien2007VIF}. (iv)~\emph{Partial
correlation with the target}: $\mathrm{pcor}(z, y \mid X)$, the
correlation of $z$ with $y$ after linearly removing $X$, which
measures residual signal the baseline does not already supply. The
\emph{combined diagnostic} treats a candidate as carrying
non-negligible independent linear signal only if it clears
\emph{both} $\max_b|r| < 0.95$ \emph{and}
$|\mathrm{pcor}(z, y\mid X)| \ge 0.10$. This is a redundancy screen,
not a certification of QSPR utility.

\section{Results}

\subsection{Baseline redundancy}

On all three regression datasets the $30$-index baseline is highly
collinear: a large fraction of the $435$ index pairs exceed
$|r| = 0.95$, and the PCA effective rank is only $3$ components at
$95\%$ explained variance and $5$ at $99\%$. This empirical rank is
far below the structural ceiling of Theorem~\ref{thm:bid},
consistent with the degree-based BID block dominating the baseline.

\subsection{Case study: the Wuzi parametric family}

The three-parameter Wuzi index family~\cite{Paper1_Wuzi} was
screened at $100$ grid points on each regression dataset
(Table~\ref{tab:wuzi_verdict_paper1}). Across all $300$ points it
clears the pairwise screen at exactly one point (FreeSolv at
$(-1,-1,1)$, $\max|r| = 0.943$ with the modified second Zagreb
index), and the maximum partial correlation with the target
anywhere on the grid is $0.036$ --- well below $0.10$ on every
dataset. By the combined diagnostic, no Wuzi parameter triple
carries non-negligible residual linear signal under the baseline.
The pairwise screen catches only one-dimensional redundancy: at
$(-1,-1,1)$ the Wuzi value is captured by a linear combination of
the baseline on all three datasets (partial correlation identically
zero), yet the pairwise $\max|r|$ dips below $0.95$ on FreeSolv
alone --- a point that ``passes'' the pairwise screen can still be
exactly redundant in the regression sense.

\begin{table}[h]
\centering
\small
\caption{Wuzi $100$-point grid screening verdict across three QSPR datasets.}
\label{tab:wuzi_verdict_paper1}
\begin{tabular}{lcccc}
\toprule
Dataset & Pass (max $\lvert r\rvert < 0.95$) & Fail & Min max $\lvert r\rvert$ & Max $\lvert\mathrm{pcor}(W, y \mid X)\rvert$\\
\midrule
ESOL          & 0 & 100 & $0.965$ & $0.028$\\
FreeSolv      & 1 &  99 & $0.943$ & $0.036$\\
Lipophilicity & 0 & 100 & $0.965$ & $0.009$\\
\bottomrule
\end{tabular}
\end{table}

\subsection{Case study: twenty candidate indices}

To test generality, twenty alternative-family candidate indices
(information-theoretic, centrality, motif, eccentricity, and
spectral-hybrid; defined in the Supporting Information) were
screened on all four benchmarks (Table~\ref{tab:novel_combined}).
The pairwise screen alone is over-permissive on every dataset ---
$75$--$85\%$ of candidates clear it --- but the combined diagnostic
leaves only $0$--$4$ candidates as plausible carriers of
independent linear signal. The strongest passer is a degree-entropy
index on BBBP at $|\mathrm{pcor}| = 0.225$; no candidate reaches
$|\mathrm{pcor}| > 0.2$ on any other dataset. These passers should
be read as ``residual-signal positive on that target under the
baseline'', not as validated QSPR descriptors.

\begin{table}[h]
\centering
\caption{Combined screening verdict on the $20$ candidates against the $30$-index baseline on the four MoleculeNet benchmarks. $\mathrm{PASS}_{\mathrm{pair}}$: $\max|r|<0.95$; $\mathrm{PASS}_{\mathrm{combined}}$: also $|\mathrm{pcor}|\ge 0.10$.}
\label{tab:novel_combined}
\begin{tabular}{lcccr}
\toprule
Dataset & $n$ & $\mathrm{PASS}_{\mathrm{pair}}$ / $20$ & $\mathrm{PASS}_{\mathrm{combined}}$ / $20$ & $\max|\mathrm{pcor}|$\\
\midrule
ESOL          & $1\,127$ & $15$ & $4$ & $0.154$\\
FreeSolv      & $\phantom{0}639$  & $15$ & $1$ & $0.129$\\
Lipophilicity & $4\,200$ & $16$ & $0$ & $0.077$\\
BBBP$^\ddag$  & $2\,039$ & $17$ & $2$ & $0.225$\\
\bottomrule
\end{tabular}

\vspace{0.4em}
\noindent\small{Combined passers: ESOL --- InfoH\_deg $(0.154)$, CutVerts $(0.118)$, FourCyc $(0.112)$, MeanEcc $(0.109)$; FreeSolv --- Triangles $(0.107)$; Lipophilicity --- none; BBBP --- InfoH\_deg $(0.225)$, InfoH\_dist $(\approx 0.16)$. $^\ddag$ BBBP partial correlation computed via a linear probability model.}
\end{table}

\subsection{Predictive performance under screening}

Does pruning by the screen cost accuracy? A $5$-fold
cross-validated RandomForest benchmark
(Table~\ref{tab:ml_benchmark}) shows the pairwise-pruned feature
set matches the full $30$-index baseline within CV noise (mean RMSE
differences $|\Delta|\le 0.005$ on all three regression datasets,
and ROC-AUC within noise on BBBP) while using roughly four-fold
fewer features. A matched-model-class benchmark (same LASSO /
$\ell_1$-logistic model and the same folds; details in the
Supporting Information) shows that comparable parsimony --- about
two-fold fewer LASSO-active features --- comes at no statistically
detectable cost on FreeSolv and BBBP ($p > 0.05$, paired $t$-test)
and at a small but detectable $<3\%$ RMSE cost on ESOL and
Lipophilicity. The \emph{strict} combined screen is too aggressive
for use as a hard filter: it retained zero features in four of five
folds on Lipophilicity.

\begin{table}[h]
\centering
\small
\caption{RandomForest $5$-fold cross-validation performance under four feature configurations (RMSE for regression, ROC-AUC for BBBP; mean $\pm$ std across folds).}
\label{tab:ml_benchmark}
\begin{tabular}{lccccrr}
\toprule
Dataset & \texttt{full} & \texttt{pca\_95} & \texttt{pairwise\_pruned} & \texttt{combined\_pruned} & full $n$ & pairwise $n$\\
\midrule
ESOL          & $1.437 \pm 0.102$ & $1.455 \pm 0.096$ & $1.435 \pm 0.087$ & $1.475 \pm 0.123$ & $30$ & $7.4$\\
FreeSolv      & $3.537 \pm 0.228$ & $3.544 \pm 0.264$ & $3.540 \pm 0.224$ & $3.561 \pm 0.300$ & $30$ & $8.0$\\
Lipophilicity & $1.030 \pm 0.024$ & $1.065 \pm 0.018$ & $1.028 \pm 0.022$ & $1.287$ ($1/5$ folds) & $30$ & $8.0$\\
BBBP          & $0.846 \pm 0.006$ & $0.818 \pm 0.014$ & $0.860 \pm 0.010$ & $0.840 \pm 0.010$ & $30$ & $6.0$\\
\bottomrule
\end{tabular}

\vspace{0.4em}
\noindent\small{The \texttt{combined\_pruned} screen retained zero features in $4$ of $5$ folds on Lipophilicity; the value shown is from the single fold that retained one feature.}
\end{table}

\section{Discussion and conclusion}

Two points follow. First, the pairwise-correlation criterion that
dominates the topological-index literature is, on real QSPR data,
too permissive to establish that a new index carries independent
information: additionally requiring a non-negligible partial
correlation with the target is a minimal strengthening that most
candidates fail. Second, the Edge-Degree-Pair Basis explains
\emph{why} this redundancy is so pervasive for degree-based
indices --- on molecular graphs of maximum degree four the whole
BID family lives in a ten-dimensional space, so a thirty-index
baseline already nearly saturates it (empirical rank $3$--$5$), and
a new BID index competes for a few unused linear directions.

These results are deliberately framed as \emph{necessary but not
sufficient}. Clearing the combined screen does not certify that an
index is a useful QSPR descriptor; it indicates only non-negligible
residual linear signal under one baseline and dataset. The
diagnostics are linear, the basis theorem covers $\Delta\le 4$, and
the strict combined screen degenerates on at least one dataset
(Lipophilicity) and must not be used as a hard feature-selection
filter. Non-linear diagnostics, extensions of the basis result
beyond maximum degree four, threshold sensitivity, and a web
deployment are left open. We provide the pipeline as an
open-source, reusable pre-flight check so that any author proposing
a new scalar topological index can report, in a standard form,
whether it carries independent linear information beyond an
explicit classical baseline.

\section*{Data and code availability}

Code, datasets, and figures are publicly available at
\url{https://github.com/Singati2/topological-index-orthogonality};
the snapshot reported here is git tag
\texttt{v1.3-paper2-orthogonality}, to be archived on Zenodo with a
permanent DOI upon submission. Supporting Information (full proof of
Theorem~\ref{thm:bid}, the $30$-index baseline list, the twenty
candidate-index definitions, and the threshold-sensitivity and
feature-selection-comparison tables) accompanies this article.

\section*{Acknowledgements}

The authors thank colleagues for discussions of the screening
protocol and the candidate index set. Funding sources and
institutional support to be acknowledged in the final version.

% =====================================================================
\begin{thebibliography}{99}
% =====================================================================
\small

\bibitem{Wiener1947}
H.~Wiener,
\emph{Structural determination of paraffin boiling points},
J.~Am.~Chem.~Soc.~\textbf{69} (1947), 17--20.

\bibitem{Randic1975}
M.~Randi\'c,
\emph{On characterization of molecular branching},
J.~Am.~Chem.~Soc.~\textbf{97} (1975), 6609--6615.

\bibitem{Gutman1972}
I.~Gutman and N.~Trinajsti\'c,
\emph{Graph theory and molecular orbitals},
Chem.~Phys.~Lett.~\textbf{17} (1972), 535--538.

\bibitem{BorovicaninDasFurtulaGutman2017Zagreb}
B.~Borovi\'canin, K.~C.~Das, B.~Furtula, and I.~Gutman,
\emph{Bounds for Zagreb indices},
MATCH Commun.~Math.~Comput.~Chem.~\textbf{78} (2017), 17--100.

\bibitem{Rada2026LinearVDB}
J.~Rada,
\emph{A finite-dimensional linear framework for vertex-degree-based
topological indices},
MATCH Commun.~Math.~Comput.~Chem.~\textbf{96} (2026), 841--863.

\bibitem{Paper1_Wuzi}
A.~Natarajan, G.~Shiwakoti, M.~Arockiaraj,
\emph{The Wuzi Index Family and the Ten-Dimensional Ceiling of Molecular BID Descriptors},
in preparation, 2026. (Companion mathematical paper.)

\bibitem{MoleculeNet}
Z.~Wu, B.~Ramsundar, E.~N.~Feinberg, J.~Gomes, C.~Geniesse,
A.~S.~Pappu, K.~Leswing, V.~Pande,
\emph{MoleculeNet: A benchmark for molecular machine learning},
Chem.~Sci.~\textbf{9} (2018), 513--530.

\bibitem{Martins2012BBBP}
I.~F.~Martins, A.~L.~Teixeira, L.~Pinheiro, A.~O.~Falcao,
\emph{A Bayesian approach to in silico blood-brain barrier
penetration modeling},
J.~Chem.~Inf.~Model.~\textbf{52} (2012), 1686--1697.

\bibitem{Delaney2004ESOL}
J.~S.~Delaney,
\emph{ESOL: Estimating aqueous solubility directly from molecular structure},
J.~Chem.~Inf.~Comput.~Sci.~\textbf{44} (2004), 1000--1005.

\bibitem{Mobley2014FreeSolv}
D.~L.~Mobley and J.~P.~Guthrie,
\emph{FreeSolv: A database of experimental and calculated hydration free energies},
J.~Comput.~Aided Mol.~Des.~\textbf{28} (2014), 711--720.

\bibitem{LipophilicityChEMBL}
AstraZeneca, \emph{Experimental in vitro DMPK and physicochemical data
on a set of publicly disclosed compounds},
ChEMBL document CHEMBL3301361 (2015);
distributed via MoleculeNet~\cite{MoleculeNet}.

\bibitem{RDKit}
G.~Landrum et al.,
\emph{RDKit: Open-source cheminformatics},
\url{https://www.rdkit.org}.

\bibitem{TodeschiniConsonniHandbook}
R.~Todeschini and V.~Consonni,
\emph{Molecular Descriptors for Chemoinformatics}, 2nd ed.,
Wiley-VCH, Weinheim, 2009.

\bibitem{WoldEsbensenGeladi1987PCA}
S.~Wold, K.~Esbensen, and P.~Geladi,
\emph{Principal component analysis},
Chemometr.~Intell.~Lab.~Syst.~\textbf{2}(1--3) (1987), 37--52.

\bibitem{OBrien2007VIF}
R.~M.~O'Brien,
\emph{A caution regarding rules of thumb for variance inflation factors},
Quality \& Quantity \textbf{41}(5) (2007), 673--690.

\bibitem{Pedregosa2011sklearn}
F.~Pedregosa et al.,
\emph{Scikit-learn: Machine learning in Python},
J.~Mach.~Learn.~Res.~\textbf{12} (2011), 2825--2830.

\bibitem{Gutman2013Degree}
I.~Gutman,
\emph{Degree-based topological indices},
Croat.~Chem.~Acta \textbf{86}(4) (2013), 351--361.

\end{thebibliography}
\end{document}
